\documentclass[11pt,a4paper]{article}
\usepackage[utf8]{inputenc}
\usepackage[T1]{fontenc}
\usepackage[ngerman]{babel}
\usepackage{lmodern}
\usepackage[margin=2.3cm]{geometry}
\usepackage{graphicx}
\usepackage{enumitem}
\usepackage{booktabs}
\usepackage{xcolor}
\usepackage{tcolorbox}
\usepackage{hyperref}
\hypersetup{colorlinks=true,linkcolor=black,urlcolor=blue!60!black}
\setlength{\parindent}{0pt}
\setlength{\parskip}{6pt}
\newcommand{\ui}[1]{\textbf{#1}}
\newcommand{\bild}[2]{\includegraphics[width=#1\textwidth,frame]{../screenshots/#2}}
\usepackage[export]{adjustbox}

\title{\Huge Namio\\[6pt]\large Namen der Schulklasse lernen -- Handbuch}
\author{Version 0.21}
\date{August 2026}

\begin{document}
\maketitle

\begin{tcolorbox}[colback=red!5,colframe=red!60!black,title=Bitte zuerst lesen]
Namio speichert Namen und Fotos von Kindern. Ob und wie du solche Daten erheben und auf einem
privaten Gerät nutzen darfst, richtet sich nach DSGVO, Schulgesetz und den Vorgaben deiner Schule
(z.\,B. Einwilligung der Eltern, dienstliches Gerät). \textbf{Das musst du selbst klären.}
Namio ist ein privates Projekt ohne Gewähr; die Verantwortung für Nutzung, Gerätesperre, Sicherung
und Löschung liegt bei dir. Die App speichert alles nur lokal und verschlüsselt und sendet nichts.
\end{tcolorbox}

\tableofcontents
\newpage

\section{Was Namio macht}
Namio hilft dir, die Namen deiner Klassen zu lernen -- und zwar so, wie du sie im Klassenraum brauchst:
Gesicht sehen, Namen wissen. Dafür gibt es fünf Quizformen, ein Lernsystem mit Wiederholung in kurzen
Abständen (Leitner-Boxen) und einen Sitzplan, der selbst als Quiz dient.

Grundsätze:
\begin{itemize}[nosep]
  \item \textbf{Kein Internet.} Die App hat keine Netzwerkberechtigung. Nichts verlässt das Gerät.
  \item \textbf{Alles verschlüsselt.} Datenbank mit SQLCipher, Schlüssel im Android-Keystore, Fotos nur im
        App-internen Speicher (sie tauchen nie in der Galerie auf).
  \item \textbf{App-Sperre.} Beim Öffnen per Fingerabdruck, Gesicht oder Geräte-PIN.
  \item \textbf{Echtes Löschen.} Klasse weg = Fotodateien weg.
\end{itemize}

\section{Installation und erster Start}
\begin{enumerate}[nosep]
  \item APK aus dem Store installieren (Obtainium erkennt Updates automatisch).
  \item Beim ersten Start: Sperre bestätigen, Hinweis lesen und mit \ui{Verstanden} bestätigen.
  \item Die Demoklasse \ui{7b} (24 Kinder mit Avataren, fertiger Sitzplan) ist zum Ausprobieren da.
        Du kannst sie jederzeit löschen (\ui{Menü (drei Punkte)} an der Klasse); sie kommt nicht zurück.
\end{enumerate}

\section{Klassen}
\begin{center}\bild{0.28}{klassen.png}\hspace{1em}\bild{0.28}{klasse.png}\end{center}
\begin{itemize}[nosep]
  \item \ui{+} legt eine Klasse an (Name, Schule, Schuljahr). Die Karte zeigt Schülerzahl und Lernfortschritt.
  \item \ui{Menü (drei Punkte)} an der Karte: Klasse löschen -- samt Schülern, Fotos, Lernständen und Sitzplänen.
  \item In der Klasse oben: \ui{CSV-Import}, \ui{Statistik}, \ui{Sitzplan}, \ui{Quiz}.
\end{itemize}

\section{Schüler und Fotos}
\begin{center}\bild{0.28}{schueler.png}\hspace{1em}\bild{0.28}{avatare.png}\end{center}
\ui{+} in der Klasse legt ein Kind an (Vorname, Nachname, Mädchen/Junge). Das Geschlecht steuert,
welche Namen im Quiz als Ablenker taugen. Im Schülerdetail:
\begin{itemize}[nosep]
  \item \ui{Fotorunde} (Kamerasymbol in der Klasse): alle Kinder ohne Foto nacheinander, Kamera bleibt offen.
  \item \ui{Foto aufnehmen}: Sucher mit quadratischem Rahmen und Gesichtsoval, Frontkamera per Symbol.
        Das Bild wird mittig quadratisch zugeschnitten und auf 800\,px verkleinert.
  \item \ui{Aus Galerie}: Bild über den System-Bildwähler -- Namio bekommt nur dieses eine Bild.
  \item \ui{Avatar wählen}: 35 Cartoon-Gesichter für Tests oder wenn kein Foto erlaubt ist.
  \item \ui{Spitzname} wird, wenn gesetzt, im Raster und Sitzplan angezeigt und gilt beim Tippen als richtig.
  \item Unten: \ui{Lernstand} pro Quizmodus (Box 1--5).
\end{itemize}

\subsection{CSV-Import}
28 Namen abtippen will niemand. \ui{CSV-Import} (Symbol in der Klasse) liest Listen aus Untis oder Excel:
Trennzeichen (Semikolon, Komma, Tab) und Zeichensatz (UTF-8 mit/ohne BOM, Windows-1252) werden
erkannt, die Spaltenzuordnung (Vorname, Nachname, „Nachname, Vorname“ in einer Spalte, Geschlecht
m/w) bestätigst du mit Vorschau. Kinder ohne Geschlechtsangabe bekommen die gewählte Vorgabe.

\section{Quiz}
\begin{center}\bild{0.28}{auswahl.png}\hspace{1em}\bild{0.28}{frage.png}\hspace{1em}\bild{0.28}{ergebnis.png}\end{center}
\begin{tabular}{@{}p{4.2cm}p{10.5cm}@{}}
\toprule
Modus & Was passiert \\
\midrule
Foto $\to$ Name (Auswahl) & Foto oben, vier Namen. Der Einstieg. \\
Foto $\to$ Name (Tippen) & Name eintippen. Tippfehler und fehlende Akzente werden toleriert (1 Fehler bei kurzen, 2 bei längeren Namen). Nur der Vorname zählt -- außer die Klasse hat zwei gleiche Vornamen. \\
Name $\to$ Foto & Name oben, neun Gesichter. \\
Sitzplan & „Wo sitzt …?“ -- auf den Platz tippen. Trainiert genau die Situation im Raum. \\
Speedrun & 60 Sekunden, so viele wie möglich. Ohne Einfluss auf den Lernstand. \\
\bottomrule
\end{tabular}

Ablenker (falsche Antworten) sind nie zufällig beliebig: zuerst Kinder, mit denen du das Ziel schon
verwechselt hast, dann gleiche Anfangsbuchstaben, dann Zufall -- immer gleiches Geschlecht, immer
dieselbe Klasse.

\subsection{Das Lernsystem}
Jedes Kind hat pro Modus eine Karte in einer von fünf Boxen. Richtig = eine Box höher, falsch = zurück
auf Box~1 (und in derselben Runde noch einmal). Fälligkeiten: Box~1 sofort, Box~2 nach 10 Minuten,
Box~3 nach 1 Tag, Box~4 nach 3 Tagen, Box~5 nach 7 Tagen. Eine Runde nimmt erst alle fälligen Karten,
dann bis zu fünf neue (einstellbar). Ist nichts fällig, kannst du trotzdem üben.

Nach der Runde: Trefferquote, Fehlerliste mit „verwechselt mit“ und \ui{Fehler wiederholen}.

\subsection{Statistik}
Pro Klasse und Modus: Verteilung der Boxen, Trefferquote der letzten Runden und die häufigsten
Verwechslungspaare -- die Kinder, die du dir gezielt ansehen solltest.

\section{Sitzplan}
\begin{center}\bild{0.28}{sitzplan.png}\hspace{1em}\bild{0.28}{sitzplan_quiz.png}\end{center}
Der Sitzplan ist frei: Tische mit 1--4 Plätzen, beliebig platziert und gedreht, Möbel (Pult, PC) als
Beschriftung. Ein Plan startet aus einer Vorlage (Doppeltischreihen, U-Form, Gruppentische, je Kind
ein Tisch) oder leer; auf Wunsch gleich mit allen Kindern der Reihe nach besetzt. Mehrere Pläne pro
Klasse (Klassenraum, Fachraum, Prüfung), einer ist Standard und dient dem Quiz.

\subsection{Ansehen (gesperrt)}
Pläne mit Tischen öffnen \textbf{gesperrt} -- nichts verrutscht aus Versehen. Verfügbar sind:
\ui{Würfel} (zufällige Person, ihr Platz blinkt), \ui{Blickrichtung} (vom Pult aus oder wie auf Papier),
\ui{Mischen} (nach Rückfrage), Zoom (zwei Finger oder Lupen), \ui{Menü (drei Punkte)} $\to$ \ui{Als PDF speichern}.

\subsection{Bearbeiten (Schloss antippen)}
\begin{itemize}[nosep]
  \item \textbf{Kind setzen/tauschen}: Kind in der Leiste oder auf einem Platz antippen, dann Ziel antippen.
        Belegtes Ziel = Tausch. Oder ziehen.
  \item \textbf{Tisch markieren}: Tisch antippen (bei belegtem Platz lange drücken). Werkzeugleiste:
        $\pm$15°, $\pm$45°, Plätze +/$-$, Breite +/$-$, Duplizieren, Beschriften, Löschen.
  \item \textbf{Mehrere Tische}: nacheinander markieren; gemeinsam ziehen, drehen (Buttons oder
        Zwei-Finger-Drehung), ausrichten (Zeile, Spalte, Drehung, Abstand).
  \item \textbf{Neuer Tisch}: Art unten wählen (Einzel, Zweier mit einem Kind, Doppel, Dreier), freie Fläche antippen.
  \item \ui{Rückgängig} für die letzten 30 Schritte. \ui{Menü (drei Punkte)}: neuer Plan, Name/Raum/Einrasten, Plan kopieren, als Standard, löschen.
\end{itemize}
Tipp: Der Editor ist Tablet-Arbeit. Am Handy sind Ansehen, Auslosen, Mischen und das Sitzplan-Quiz
bequem; der Plan ist dort automatisch so gezoomt, dass man scrollt statt zu schrumpfen.

\section{Einstellungen und Daten}
\begin{itemize}[nosep]
  \item \ui{App-Sperre} an/aus (Standard an). Ohne Gerätesperre bleibt die App offen.
  \item \ui{Sitzplan-Blickrichtung}, \ui{Neue Karten pro Runde} (1--15).
  \item \ui{Exportieren}: alle Daten als Datei \texttt{namio-Datum.namio}, verschlüsselt mit einem Passwort
        (AES-256, Passwort wird nirgends gespeichert). \ui{Importieren} ersetzt alle Daten auf dem Gerät --
        so ziehst du vom Tablet aufs Handy um oder sicherst dich ab.
  \item \ui{Alle Daten löschen}: Klassen, Fotos, Lernstände, Sitzpläne -- endgültig.
  \item Zwischen Mitte Juni und Ende August erinnert Namio einmal, abgegebene Klassen zu löschen.
\end{itemize}

\section{Tipps aus der Praxis}
\begin{itemize}[nosep]
  \item Erst Fotos, dann Quiz: Ohne Foto kann ein Kind nicht abgefragt werden.
  \item Montagmorgen: fünf Minuten \emph{Foto $\to$ Name (Auswahl)}, dann \emph{Sitzplan}.
  \item Verwechslungspaare in der Statistik sind die beste Lernliste.
  \item Am Schuljahresende exportieren (falls du etwas behalten willst) und dann löschen.
\end{itemize}

\section{Häufige Fragen}
\textbf{Wo liegen die Fotos?} Nur im internen App-Speicher, verschlüsselt erreichbar nur für Namio. Nie in der Galerie.

\textbf{Kann ich Fotos für alle Kinder auf einmal machen?} Ja: In der Klasse das Kamerasymbol antippen
(erscheint, solange Kinder ohne Foto da sind). Die Kamera bleibt offen, oben steht das nächste Kind,
jedes Foto springt automatisch weiter; \ui{Überspringen} für Kinder, die gerade fehlen.

\textbf{Warum zeigt die Quizauswahl „Nichts fällig“?} Alle Karten sind in höheren Boxen. Üben geht trotzdem;
fällig werden sie nach den Intervallen oben.

\textbf{Was passiert bei Wechsel des Geräts?} Exportieren, Datei rüberkopieren (Dateimanager, NAS, USB), importieren.

\textbf{Ist die App barrierefrei?} Bedienbar mit TalkBack und großer Schrift; Fotos werden mit Namen vorgelesen.

\end{document}
